\subsection{Repository Structure}\label{repository-structure}

Github was used as our team's code repository management system, and
from the root of
\href{https://github.com/Eines-Informatiques-Avancades/Project_2026_II}{the
project}, the code is divided among 3 main folders: \texttt{src},
\texttt{results}, and \texttt{docs}.

\begin{itemize}
\tightlist
\item
  \texttt{src} contains:

  \begin{itemize}
  \tightlist
  \item
    Makefile \& shell scripts
  \item
    Various versions of the main program (serial \& parallel)
  \item
    A \texttt{confs} folder where initial configurations are stored and
    compile-time options can be controlled via a file called
    \texttt{input.dat}
  \item
    A \texttt{lib} folder where module source files, python analysis
    files, and plotting style dependencies reside (\emph{further detail
    below})
  \end{itemize}
\item
  \texttt{results} contains the output data from the simulations
\item
  \texttt{docs} contains documentation and reports, including
  \texttt{img} folder for embedded images
\end{itemize}

In addition to these folders, a customary \texttt{README} file is
included in the root directory detailing run instructions and code
dependencies/requirements. The bulk of the code is written in Fortran,
while the post-processing is done in Python.

\#\#\# Module Library Introduction The following modules support the
\emph{\texttt{main...f90}} files: - \texttt{energy.f90} \&
\texttt{energy\_all\_atoms.f90} - Permit \textbf{dual energy modes},
computing energetic contributions using either TraPPE-UA or OPLS-AA
(with effective backbone potentials) depending on the
\texttt{explicit\_h} toggle dynamically read from \texttt{input.dat}. -
\texttt{initial-conf.f90} - Generates initial configuration of the
polymer chain. - \texttt{io.f90} - File input/ouput module. -
\texttt{monte\_carlo.f90} - Provides the rotation algorithm and
subroutine MCS step used to update spatial configurations. -
\texttt{observables.f90} - Computes the structural observables:
End-to-end distance (\(R_{ee}\)), Radius of Gyration (\(R_g\)), and
torsion angles. - \texttt{parameters.f90} - Stores global parameters
abstracted from \emph{\texttt{main...f90}} code. Works in tandem with
\texttt{input.dat} settings. - \emph{\texttt{plot\_...py}} files -
Various scripts used for plotting observables and serial vs parallel
code comparisons. - \texttt{requirements.txt} - List of python module
dependencies. - \texttt{science.mplstyle} - stylized document used in
plotting to control visualization theme.

\subsection{Code Management via
Github}\label{code-management-via-github}

The project leader was responsible for reviewing all pull requests,
ensuring that the code followed the project guidelines, and merging the
code into the main branch. The team coordinated through a WhatsApp group
chat and would regularly check in on each other's progress. Everyone was
informed when code was merged so we could all pull the latest changes.

Most times, approving and merging was straightforward with no issues,
but there was an interesting challenge when multiple uploads from
different team members overlapped, causing conflicts. This forced us to
learn about stashing and rebasing to the newly updated master branch
before recommitting. The following process was ironed out to ensure a
smooth integration:

\begin{enumerate}
\def\labelenumi{\arabic{enumi}.}
\item
  Stash local (uncommitted) changes

\begin{Shaded}
\begin{Highlighting}[]
   \FunctionTok{git}\NormalTok{ stash}
\end{Highlighting}
\end{Shaded}
\item
  Pull latest changes from master branch

\begin{Shaded}
\begin{Highlighting}[]
   \FunctionTok{git}\NormalTok{ checkout master}
   \FunctionTok{git}\NormalTok{ pull main master}
\end{Highlighting}
\end{Shaded}
\item
  Rebase local changes onto master branch so it starts after the new
  merged code

\begin{Shaded}
\begin{Highlighting}[]
   \FunctionTok{git}\NormalTok{ checkout }\OperatorTok{\textless{}}\NormalTok{fork{-}feature}\OperatorTok{\textgreater{}}\NormalTok{/}\OperatorTok{\textless{}}\NormalTok{fork{-}branch}\OperatorTok{\textgreater{}}
   \FunctionTok{git}\NormalTok{ rebase master}
\end{Highlighting}
\end{Shaded}
\item
  Resolve any conflicts \& re-apply stashed changes

\begin{Shaded}
\begin{Highlighting}[]
   \FunctionTok{git}\NormalTok{ stash pop}
\end{Highlighting}
\end{Shaded}
\item
  Commit and push (using \texttt{-\/-force} parameter due to rebase)

\begin{Shaded}
\begin{Highlighting}[]
   \FunctionTok{git}\NormalTok{ push }\OperatorTok{\textless{}}\NormalTok{fork{-}remote}\OperatorTok{\textgreater{}} \OperatorTok{\textless{}}\NormalTok{fork{-}branch}\OperatorTok{\textgreater{}}\NormalTok{ {-}{-}force}
\end{Highlighting}
\end{Shaded}
\end{enumerate}

\subsection{Makefile}\label{makefile}

As the project evolved, so did the Makefile. Here are some of the
features we implemented:

Starting off, there are few main variables that control the compilation
process.... \begin{Shaded}
\begin{Highlighting}[]
\DataTypeTok{FC}       \CharTok{=}\StringTok{ gfortran}
\DataTypeTok{FFLAGS}   \CharTok{=}\StringTok{ {-}O2 {-}Wall {-}Wextra}
\DataTypeTok{OBJ\_DIR}  \CharTok{=}\StringTok{ ../bin}
\DataTypeTok{EXE}      \CharTok{=}\StringTok{ }\CharTok{$(}\DataTypeTok{OBJ\_DIR}\CharTok{)}\StringTok{/main\_serial.x}
\DataTypeTok{LIB\_OBJ}  \CharTok{=}\StringTok{ }\CharTok{$(}\DataTypeTok{OBJ\_DIR}\CharTok{)}\StringTok{/parameters.o }\CharTok{\textbackslash{}}
\StringTok{           }\CharTok{$(}\DataTypeTok{OBJ\_DIR}\CharTok{)}\StringTok{/io.o }\CharTok{\textbackslash{}}
\StringTok{           ...}
\DataTypeTok{MAIN\_OBJ} \CharTok{=}\StringTok{ }\CharTok{$(}\DataTypeTok{OBJ\_DIR}\CharTok{)}\StringTok{/main\_serial.o}

\DataTypeTok{NP}          \CharTok{?=}\StringTok{ 4}
\DataTypeTok{OMP\_THREADS} \CharTok{?=}\StringTok{ 2}
\end{Highlighting}
\end{Shaded}

\begin{itemize}
\tightlist
\item
  \texttt{FC}: Stands for ``Fortran Compiler''. We tell it to use
  gfortran.
\item
  \texttt{FFLAGS}: These are the compiler flags. \texttt{-O2} is the
  optimization schema, and \texttt{-Wall\ -Wextra} enables warnings.
\item
  \texttt{OBJ\_DIR} \& \texttt{EXE}: Defines the compiled binaries
  directory and the name of the final executable.
\item
  \texttt{LIB\_OBJ} \& \texttt{MAIN\_OBJ}: Defines the object files for
  the library and the main program.
\item
  \texttt{NP} \& \texttt{OMP\_THREADS}: Defines the number of processes
  (MPI) and threads (openMP) to use for the parallel version of the
  program. The default values are set to 4 and 2, respectively, but they
  can be overridden by setting them when compiling in the command line
  (e.g., \texttt{make\ run\_parallel\ NP=7\ OMP\_THREADS=3}).
\end{itemize}

In order to encompass both the serial and parallel versions of the
program, a check is performed on the \texttt{make} command entered on
the command-line for the word \texttt{parallel}. If found, the user's
system is searched to ensure the environment has the necessary MPI and
OpenMP libraries installed. If not, an error message is printed and the
program exits. If they are installed, the compiler variable \texttt{FC}
is overwritten and the MPI and OpenMP flags are added to the
\texttt{FFLAGS} variable.

\begin{Shaded}
\begin{Highlighting}[]
\ControlFlowTok{ifneq}\NormalTok{ (,}\CharTok{$(}\KeywordTok{findstring}\StringTok{ parallel}\KeywordTok{,}\CharTok{$(}\DataTypeTok{MAKECMDGOALS}\CharTok{))}\NormalTok{)}
  \CommentTok{\# MPI Compiler Wrapper Verification}
  \DataTypeTok{MPI\_BIN} \CharTok{:=}\StringTok{ }\CharTok{$(}\KeywordTok{shell}\StringTok{ command {-}v mpif90 2\textgreater{}/dev/null}\CharTok{)}
  \ControlFlowTok{ifeq}\NormalTok{ (}\CharTok{$(}\KeywordTok{strip}\StringTok{ }\CharTok{$(}\DataTypeTok{MPI\_BIN}\CharTok{))}\NormalTok{,)}
    \CharTok{$(}\KeywordTok{error}\StringTok{ "MPI Compiler is required but \textquotesingle{}mpif90\textquotesingle{} was not found! Please install it (e.g. \textquotesingle{}sudo apt install openmpi{-}bin\textquotesingle{})"}\CharTok{)}
  \ControlFlowTok{endif}
  \CommentTok{\# Overwrite Compiler for MPI parallelization flag linking}
  \DataTypeTok{FC} \CharTok{=}\StringTok{ mpif90}
\NormalTok{  ...}
\end{Highlighting}
\end{Shaded}

Instead of writing a rule for every single \texttt{.f90} file, a pattern
rule (with \texttt{\%}) is used.

\begin{Shaded}
\begin{Highlighting}[]
\DecValTok{$(OBJ\_DIR)/\%.o:}\DataTypeTok{ lib/\%.f90}
\ErrorTok{    }\CharTok{@}\FunctionTok{mkdir {-}p }\CharTok{$(}\DataTypeTok{OBJ\_DIR}\CharTok{)}
    \CharTok{$(}\DataTypeTok{FC}\CharTok{)} \CharTok{$(}\DataTypeTok{FFLAGS}\CharTok{)}\NormalTok{ {-}J}\CharTok{$(}\DataTypeTok{OBJ\_DIR}\CharTok{)}\NormalTok{ {-}I}\CharTok{$(}\DataTypeTok{OBJ\_DIR}\CharTok{)}\NormalTok{ {-}c }\CharTok{$\textless{}}\NormalTok{ {-}o }\CharTok{$@}
\end{Highlighting}
\end{Shaded}

\begin{itemize}
\tightlist
\item
  It says: ``To make any \texttt{.o} file in the object directory
  (\texttt{bin/}) from a matching \texttt{.f90} file in \texttt{lib/},
  run this command.''
\item
  \texttt{-c\ \$\textless{}}: Compiles the ``first dependency''
  (\texttt{\$\textless{}}, which is the \texttt{.f90} file) without
  linking (because that's done in the \texttt{\$(EXE)} target).
\item
  \texttt{-J\$(OBJ\_DIR)\ -I\$(OBJ\_DIR)}: Tells Fortran to put module
  \texttt{.mod} files in the \texttt{bin/} folder.
\end{itemize}

The following is an example of the \texttt{all} target (serial version
of the program) with its executable structure, which we then follow up
with the explicit dependencies to create the \texttt{.o} files for the
library modules and the main program

\begin{Shaded}
\begin{Highlighting}[]
\DecValTok{all:}\DataTypeTok{ }\CharTok{$(}\DataTypeTok{EXE}\CharTok{)}

\DecValTok{$(EXE):}\DataTypeTok{ }\CharTok{$(}\DataTypeTok{LIB\_OBJ}\CharTok{)}\DataTypeTok{ }\CharTok{$(}\DataTypeTok{MAIN\_OBJ}\CharTok{)}
\ErrorTok{    }\CharTok{@}\FunctionTok{mkdir {-}p }\CharTok{$(}\DataTypeTok{OBJ\_DIR}\CharTok{)}
    \CharTok{$(}\DataTypeTok{FC}\CharTok{)} \CharTok{$(}\DataTypeTok{FFLAGS}\CharTok{)}\NormalTok{ {-}o }\CharTok{$@} \CharTok{$(}\DataTypeTok{LIB\_OBJ}\CharTok{)} \CharTok{$(}\DataTypeTok{MAIN\_OBJ}\CharTok{)}

\DecValTok{$(OBJ\_DIR)/energy.o:}\DataTypeTok{ lib/energy.f90 }\CharTok{\textbackslash{}}
\DataTypeTok{                     }\CharTok{$(}\DataTypeTok{OBJ\_DIR}\CharTok{)}\DataTypeTok{/parameters.o}
\NormalTok{...}

\DecValTok{$(OBJ\_DIR)/main\_serial.o:}\DataTypeTok{ main\_serial.f90 }\CharTok{$(}\DataTypeTok{LIB\_OBJ}\CharTok{)}
\ErrorTok{    }\CharTok{@}\FunctionTok{mkdir {-}p }\CharTok{$(}\DataTypeTok{OBJ\_DIR}\CharTok{)}
    \CharTok{$(}\DataTypeTok{FC}\CharTok{)} \CharTok{$(}\DataTypeTok{FFLAGS}\CharTok{)}\NormalTok{ {-}J}\CharTok{$(}\DataTypeTok{OBJ\_DIR}\CharTok{)}\NormalTok{ {-}I}\CharTok{$(}\DataTypeTok{OBJ\_DIR}\CharTok{)}\NormalTok{ {-}c }\CharTok{$\textless{}}\NormalTok{ {-}o }\CharTok{$@}
\end{Highlighting}
\end{Shaded}... \begin{itemize}
\tightlist
\item
  \texttt{all}: This is the default target that needs \texttt{\$(EXE)}
  to exist.
\item
  \texttt{\$(EXE)}: To build the final executable, it needs all the .o
  files from the \texttt{../bin} folder, which it makes sure exists
  (\texttt{mkdir\ -p}), and then links them all together using
  \texttt{gfortran} or \texttt{mpif90} to output (\texttt{-o\ \$@}) the
  final program.

  \begin{itemize}
  \tightlist
  \item
    \textbf{Note}: The \texttt{@} symbol is used 2 different ways here:

    \begin{enumerate}
    \def\labelenumi{\arabic{enumi}.}
    \tightlist
    \item
      Before the \texttt{mkdir} command to suppress being printed to the
      terminal.
    \item
      As an automatic variable representing the target name
      (\texttt{\$(EXE)} in this case).
    \end{enumerate}
  \end{itemize}
\end{itemize}

Last, we declare the shortcut commands as \texttt{.PHONY} targets to
prevent conflicts with files that may have the same name as a target.

\begin{Shaded}
\begin{Highlighting}[]
\OtherTok{.PHONY:}\DataTypeTok{ all clean run figures pipeline}

\DecValTok{run:}\DataTypeTok{ all}
\ErrorTok{    }\CharTok{@}\FunctionTok{mkdir {-}p ../results}
    \CharTok{$(}\DataTypeTok{EXE}\CharTok{)}

\DecValTok{figures:}
\ErrorTok{    }\NormalTok{python lib/plot\_results.py}

\DecValTok{clean:}\DataTypeTok{ }
\ErrorTok{    }\NormalTok{rm {-}rf }\CharTok{$(}\DataTypeTok{OBJ\_DIR}\CharTok{)}\NormalTok{/*.o }\CharTok{$(}\DataTypeTok{OBJ\_DIR}\CharTok{)}\NormalTok{/*.mod }\CharTok{$(}\DataTypeTok{EXE}\CharTok{)}

\DecValTok{pipeline:}
\ErrorTok{    }\CharTok{$(}\DataTypeTok{MAKE}\CharTok{)}\NormalTok{ run}
    \CharTok{$(}\DataTypeTok{MAKE}\CharTok{)}\NormalTok{ figures}
\end{Highlighting}
\end{Shaded}

\begin{itemize}
\tightlist
\item
  This enables the following compile-time shortcuts:

  \begin{itemize}
  \tightlist
  \item
    \texttt{make\ run}: Automatically builds the code (\texttt{all}) and
    then executes it.
  \item
    \texttt{make\ figures}: Envokes python to generate the plots from
    the \texttt{../results} folder.
  \item
    \texttt{make\ clean}: Acts like a reset button, deleting all
    compiled files to start fresh.
  \item
    \texttt{make\ pipeline}: A neat wrapper we built to run the code and
    generate the python plots back-to-back.
  \end{itemize}
\end{itemize}

\subsection{Shell Scripting}\label{shell-scripting}
% Itxaso bit added + Arthur's explanation merged

Shell scripting was used to make our runs compatible with the CERQT2
cluster environment and to automate benchmarking. The scripts
initialize the runtime environment, load the required modules, and set
variables such as \texttt{OMP\_NUM\_THREADS}, \texttt{NSLOTS}, and
\texttt{MPLBACKEND=Agg} for stable non-interactive execution. They were
also used to automate parameter sweeps, build explicit trajectory
manifests for the observables analysis, and collect timing data into
summary files for MPI and OpenMP tests.

A qsub script \texttt{run\_collab.qsub} was also written during the
testing of the program to submit an individual portion of the simulation
to the HPC cluster. It involved loading the appropriate modules for that
queue, selecting a custom makefile (separate from the main program's
Makefile), indentifying the correct commandline arguments to pass at
compile time, and submit the mpirun command.